\documentclass[11pt,a4paper]{article}
\usepackage[utf8]{inputenc}
\usepackage{amsmath,amssymb,amsthm}
\usepackage{listings}
\usepackage[margin=2.5cm]{geometry}
\usepackage{xcolor}

\newtheorem{theorem}{Theorem}
\newtheorem{lemma}[theorem]{Lemma}

\lstset{
  language=C++,
  basicstyle=\ttfamily\small,
  keywordstyle=\color{blue}\bfseries,
  commentstyle=\color{green!50!black},
  stringstyle=\color{red!70!black},
  numbers=left,
  numberstyle=\tiny\color{gray},
  breaklines=true,
  frame=single,
  tabsize=4,
  showstringspaces=false
}

\title{IOI 1993 -- Day 2, Task 1: Operations}
\author{Solution}
\date{}

\begin{document}
\maketitle

\section{Problem Statement}

Given a starting integer $s$ and a target integer $t$, find the minimum
number of operations to transform $s$ into $t$. The allowed operations are:
\begin{itemize}
  \item Add 1: $x \to x + 1$
  \item Subtract 1: $x \to x - 1$
  \item Multiply by 2: $x \to 2x$
  \item Integer divide by 2: $x \to \lfloor x/2 \rfloor$
\end{itemize}

Output the minimum number of steps and the operation sequence.

\section{Solution}

\subsection{BFS on the State Space}

This is a shortest-path problem in an unweighted graph where states are
integers and transitions are the four operations. BFS finds the minimum
number of steps.

The key question is bounding the search range. For the operation set
$\{+1, -1, \times 2, \div 2\}$:

\begin{lemma}
An optimal path from $s$ to $t$ only visits values in
$[\min(s,t) - 2,\; 2 \cdot \max(s,t) + 2]$.
\end{lemma}

This bound is loose but sufficient; the actual BFS terminates quickly because
doubling shrinks the gap fast.

\subsection{Alternative: Working Backwards}

A cleaner approach works backwards from the target:
\begin{itemize}
  \item If $t > s$ and $t$ is even, halve $t$ (inverse of $\times 2$).
  \item If $t > s$ and $t$ is odd, add 1 to $t$ first, then halve
    (inverse of $-1$ then $\times 2$).
  \item If $t \le s$, the answer is $s - t$ steps of $-1$.
\end{itemize}

This greedy backward approach gives $O(\log t)$ steps. However, it does not
always produce the optimal sequence for the full 4-operation set (the BFS
approach is exact).

\section{Complexity Analysis}

\begin{itemize}
  \item \textbf{BFS:} $O(|V| + |E|)$ where $|V|$ is bounded by the search
    range, typically $O(\max(s,t))$.
  \item \textbf{Backward greedy:} $O(\log t)$ time.
\end{itemize}

\section{Implementation: BFS}

\begin{lstlisting}
#include <iostream>
#include <queue>
#include <unordered_map>
#include <vector>
#include <algorithm>
using namespace std;

int main() {
    int start, target;
    cin >> start >> target;

    if (start == target) {
        cout << 0 << endl;
        return 0;
    }

    unordered_map<int, int> dist;
    unordered_map<int, int> parent;
    unordered_map<int, string> op;

    queue<int> q;
    q.push(start);
    dist[start] = 0;

    int lo = min(start, target) - 2;
    int hi = max(start, target) * 2 + 2;

    while (!q.empty()) {
        int cur = q.front(); q.pop();

        struct Move { int nxt; string name; };
        Move moves[] = {
            {cur + 1, "+1"},
            {cur - 1, "-1"},
            {cur * 2, "*2"},
            {cur / 2, "/2"}
        };

        for (auto& [nxt, name] : moves) {
            if (nxt < lo || nxt > hi) continue;
            if (dist.count(nxt)) continue;

            dist[nxt] = dist[cur] + 1;
            parent[nxt] = cur;
            op[nxt] = name;

            if (nxt == target) {
                cout << dist[nxt] << endl;
                vector<string> path;
                int x = target;
                while (x != start) {
                    path.push_back(op[x]);
                    x = parent[x];
                }
                reverse(path.begin(), path.end());
                for (auto& s : path)
                    cout << s << " ";
                cout << endl;
                return 0;
            }
            q.push(nxt);
        }
    }

    cout << -1 << endl; // unreachable (should not happen)
    return 0;
}
\end{lstlisting}

\section{Implementation: Backward Greedy}

This approach works backwards from the target and produces the forward
operation sequence. It is optimal for the operation set $\{+1, -1, \times 2\}$
but may not be optimal when $\div 2$ is also available.

\begin{lstlisting}
#include <iostream>
#include <vector>
#include <algorithm>
using namespace std;

int main() {
    int start, target;
    cin >> start >> target;

    vector<string> revOps; // operations in reverse order
    int cur = target;

    while (cur > start) {
        if (cur % 2 == 0) {
            cur /= 2;
            revOps.push_back("*2");
        } else {
            cur++;
            revOps.push_back("-1");
        }
    }

    // Now cur <= start; need (start - cur) subtract-1 operations
    for (int i = 0; i < start - cur; i++)
        revOps.push_back("-1");

    reverse(revOps.begin(), revOps.end());

    cout << revOps.size() << endl;
    for (auto& s : revOps)
        cout << s << " ";
    cout << endl;

    return 0;
}
\end{lstlisting}

\section{Notes}

\begin{itemize}
  \item The BFS approach guarantees optimality for any operation set.
  \item The backward greedy approach is simpler and runs in $O(\log t)$
    time, but is only provably optimal for $\{+1, -1, \times 2\}$.
  \item For the original IOI constraints, BFS within a bounded range is
    sufficient.
\end{itemize}

\end{document}
